
%allgemeine Formatangaben
\documentclass[
 a4paper, 					% Papierformat
 11pt,						% Schriftgr��e
 ngerman, 					% f�r Umlaute, Silbentrennung etc.
 titlepage,					% es wird eine Titelseite verwendet
 bibliography=totoc,		% Literaturverzeichnis im Inhaltsverzeichnis auff�hren
 listof=totoc,				% Verzeichnisse im Inhaltsverzeichnis auff�hren
 oneside, 					% einseitiges Dokument
 captions=nooneline,		% einzeilige Gleitobjekttitel ohne Sonderbehandlung wie mehrzeilige Gleitobjekttitel behandeln
 numbers=noenddot,			% �berschriften-??Nummerierung ohne Punkt am Ende
 parskip=half				% zwischen Abs�tzen wird eine halbe Zeile eingef�gt
 ]{scrreprt}

\newcommand{\titel}{Beispieldokument f�r die Erstellung einer Bachelor- oder Masterarbeit mit LaTeX}
\newcommand{\abschlussart}{Master of Science (M.Sc.)}
\newcommand{\arbeit}{Masterarbeit}
\newcommand{\hochschule}{Hochschule f�r Technik, Wirtschaft und Kultur Leipzig}
\newcommand{\fakultaet}{Fakult�t Informatik und Medien}
\newcommand{\autor}{Max Muster}
\newcommand{\studiengang}{Masterstudiengang Medieninformatik}
\newcommand{\matrikelnr}{12345}
\newcommand{\erstgutachter}{Prof. Dr. Musterfrau}
\newcommand{\zweitgutachter}{Prof. Dr. Mustermann}
\newcommand{\ort}{Leipzig}
 	% einbinden von pers�nlichen Daten

%Highlighting
%\usepackage{color}
\usepackage{soul}


% Anpassung an Landessprache
\usepackage[ngerman]{babel}
 
\usepackage[default]{sourcesanspro}
 
% Verwenden von Sonderzeichen und Silbentrennung
\usepackage[latin1]{inputenc}	
\usepackage[T1]{fontenc}			
\usepackage{textcomp} 																% Euro-Zeichen und andere
\usepackage[babel,german=quotes]{csquotes}						% Anf�hrungszeichen
\RequirePackage[ngerman=ngerman-x-latest]{hyphsubst} 	% erweiterte Silbentrennung

% Befehle aus AMSTeX f�r mathematische Symbole z.B. \boldsymbol \mathbb
\usepackage{amsmath,amsfonts}

% Zeilenabst�nde und Seitenr�nder 
\usepackage{setspace}
\usepackage{geometry}

% Einbinden von JPG-Grafiken
\usepackage{graphicx}

% zum Umflie�en von Bildern
% Verwendung unter http://de.wikibooks.org/wiki/LaTeX-Kompendium:_Baukastensystem#textumflossene_Bilder
\usepackage{floatflt}

% Verwendung von vordefinierten Farbnamen zur Colorierung
% Palette und Verwendung unter http://kitt.cl.uzh.ch/kitt/CLinZ.CH/src/Kurse/archiv/LaTeX-Kurs-Farben.pdf
\usepackage[usenames,dvipsnames]{color} 

% Tabellen
\usepackage{array}
\usepackage{longtable}

% einfache Grafiken im Code
% Einf�hrung unter http://www.math.uni-rostock.de/~dittmer/bsp/pstricks-bsp.pdf
\usepackage{pstricks}

% Quellcodeansichten
\usepackage{verbatim}
\usepackage{moreverb} 											% f�r erweiterte Optionen der verbatim Umgebung
% Befehle und Beispiele unter http://www.ctex.org/documents/packages/verbatim/moreverb.pdf
\usepackage{listings} 											% f�r angepasste Quellcodeansichten siehe
% Kurzeinf�hrung unter http://blog.robert-kummer.de/2006/04/latex-quellcode-listing.html

% Glossar und Abbildungsverzeichnis
\usepackage[
nonumberlist, %keine Seitenzahlen anzeigen
acronym,      %ein Abk�rzungsverzeichnis erstellen
toc          %Eintr�ge im Inhaltsverzeichnis
]      %im Inhaltsverzeichnis auf section-Ebene erscheinen
{glossaries}

% verlinktes und Farblich angepasstes Inhaltsverzeichnis
\usepackage[pdftex,
colorlinks=true,
linkcolor=InterneLinkfarbe,
urlcolor=ExterneLinkfarbe,
citecolor=Citefarbe]{hyperref}
\usepackage[all]{hypcap}

% URL verlinken, lange URLs umbrechen
\usepackage{url}

% sorgt daf�r, dass Leerzeichen hinter parameterlosen Makros nicht als Makroendezeichen interpretiert werden
\usepackage{xspace}

% Beschriftungen f�r Abbildungen und Tabellen
\usepackage{caption}

% Entwicklerwarnmeldungen entfernen
\usepackage{scrhack}

% Keine Worttrennungen
\usepackage[none]{hyphenat}
\sloppy

% PDF/A erzeugen
%\usepackage[a-1b]{pdfx}		% einbinden der verwendeten Latex-Pakete

\singlespacing				% einfacher Zeilenabstand
%\onehalfspacing 			% 1,5facher Zeilenabstand

\definecolor{InterneLinkfarbe}{rgb}{0.1,0.1,0.3} 	% Farbliche Absetzung von externen Links
%\definecolor{InterneLinkfarbe}{rgb}{0,0,0}
\definecolor{ExterneLinkfarbe}{rgb}{0.1,0.1,0.7}	% Farbliche Absetzung von internen Links
%\definecolor{ExterneLinkfarbe}{rgb}{0,0,0}
\definecolor{Citefarbe}{rgb}{0.0,0.5,0.5}
%\definecolor{Citefarbe}{rgb}{0,0,0}

% Einstellungen f�r Fu�noten:
\captionsetup{font=footnotesize,labelfont=sc,singlelinecheck=true,margin={5mm,5mm}}

% Stil der Quellenangabe
\bibliographystyle{alphadin}
%\bibliographystyle{ACM-Reference-Format}

%Ausschluss von Schusterjungen
\clubpenalty = 10000
%Ausschluss von Hurenkindern
\widowpenalty = 10000

% Befehle, die Umlaute ausgeben, f�hren zu Fehlern, wenn sie hyperref als Optionen �bergeben werden
\hypersetup{
%    pdftitle={\titel \untertitel},
%    pdfauthor={\autor},
%    pdfcreator={\autor},
%    pdfsubject={\titel \untertitel},
%    pdfkeywords={\titel \untertitel},
}

% Beispiel f�r eine Listings-Codeumbebungen
% Bei mehreren Definitionen empfielt sich das auslagern in eine externe Datei
\lstloadlanguages{Java,HTML}
\lstset{
	frame=tb,
	framesep=5pt,
	basicstyle=\footnotesize\ttfamily,
	showstringspaces=false,
	keywordstyle=\ttfamily\bfseries\color{CadetBlue},
	identifierstyle=\ttfamily,
	stringstyle=\ttfamily\color{OliveGreen},
	commentstyle=\color{GrayBlue},
	rulecolor=\color{Gray},
	xleftmargin=5pt,
	xrightmargin=5pt,
	aboveskip=\bigskipamount,
	belowskip=\bigskipamount
} 

%Den Punkt am Ende jeder Beschreibung deaktivieren
\renewcommand*{\glspostdescription}{}

% Empfehlung: Abkuerzungsverzeichnis und Glossar sind in Graduierunsarbeiten
% nicht zwingend notwendig

% %Glossar-Befehle anschalten
% \makeglossaries
% \glsenablehyper
% \input{Header/Abkuerzungen}
% \input{Header/Glossar}


% --------- Seitenlayout -----------

% Zwischenabgabe
\geometry{
	top=2cm,
	left=3cm,
	right=3cm
}

% Endabgabe
%\areaset[0.5cm]{16cm}{24cm}


\begin{document}

%\include{Deckblatt/titel}						% Deckblatt der vorliegenden Arbeit
% Die Daten werden in Header/Metadaten.tex eingstellt
\begin{titlepage}
\begin{Large}
\begin{center}

\includegraphics[width=0.5\textwidth]{Abbildungen/HTWK_Zusatz_de_V_Black_sRGB.png}
% \textbf{\hochschule\ \ort}\\[5pt]
% \fachbereich\\
% \studiengang\\
% \vskip 1cm
% \arbeit\\
% zur Erlangung der akademischen Grades
\\[8pt]

\arbeit\\[-1mm]
zur Erlangung des akademischen Grades\\[3mm]
\abschlussart\\[3mm]
im \studiengang\\[-1mm]
der \fakultaet%\\[-1mm]
%der \hochschule  %% kann entfallen
\\
\vfill
{\LARGE\bfseries \titel \par}
\vfill

vorgelegt von\\
 \autor%\\
% Matrikelnummer: \matrikelnr  
\\[8pt]
\ort, den \today

\end{center}
\vfill
\begin{tabular}{ll}
Erstpr�fer: & \erstgutachter\\
Zweitpr�fer: & \zweitgutachter\\
\end{tabular}
\end{Large}
\end{titlepage}

\frontmatter												% Seitenzählerstart vor dem Text

\chapter*{Abstrakt}
\label{sec:Abstrakt}

Das E-Mail-Management ist Teil unseres t"aglichen Lebens. Es werden mehrere Clients genutzt, um regelm"a"sig auf die E-Mails zuzugreifen. Bei der Masse an E-Mails, die wir t"aglich Empfangen, ist es schwierig, den "Uberblick zu behalten. Eine Snooze-Funktionalit"at kann dabei helfen, das Management zu verbessern.

Das Ziel dieser Arbeit ist zu beantworten, welche M"oglichkeiten es gibt, eine Snooze-Funktionalit"at "uber das IMAP-Protokoll interoperabel zu implementieren.

Dazu wurden existierende L"osungen in E-Mail-Clients betrachtet sowie andere M"oglichkeiten untersucht, die das IMAP-Protokoll bietet. Zus"atzlich wurden diese auf die Praxistauglichkeit anhand der Nextcloud Mail-App untersucht.

Die Analyse hat gezeigt, dass es nicht die eine perfekte L"osung f"ur eine interoperable Implementierung gibt. Es wurden mehrere Teill"osungen beschrieben, die in Kombination eine ganzheitliche L"osung bieten k"onnen. Die Nachteile der einzelnen Teill"osungen wurden gegeneinander abgewogen. Als Vorschlag f"ur eine Standard-Erweiterung des IMAP-Protokolls wurde eine Kombination aus der Verwendung von Keywords, dem Nutzen von einer neuen Special-Use-Mailbox und der Verwendung der Quick-Resynchronization-Erweiterung empfohlen.  

Basierend auf diesem Vorschlag kann ein Internet-Draft entstehen und der Standadisierungsprozess f"ur eine Erweiterung beginnen.
						% Abstrakt

%% Empfehlung: Auf Danksagungen + Vorwort in wissenschaftlichen Arbeiten eher verzichten -- ggf. die entsprechenden Dateien im Ordner Inhalt anlegen und hier auskommentieren (PKW)
%% \include{Inhalt/Danksagung}					% Danksagung
%% \include{Inhalt/Vorwort}						% Vorwort

\tableofcontents										% Inhaltsverzeichnis

\mainmatter													% Seitenzählerstart Haupttext

\chapter{AT: Einf�hrendes Kapitel}
\label{sec:Intro}

\begin{itemize}
	\item Relevanz/Abgrenzung/Ziele
\end{itemize}


\begin{itemize}
	\item Zum Thema hinf�hren (Beispiele), These/Gegenthese, bedeutende Aussage durch namenhafte Pers�hnlichkeit, �berraschende Feststellung machen, (aktualit�t?) \hl{Warum ist das Thema interesant, wichtig und spannend (wissebschaftlich als auch f�r jeden der die arbeit liest)}
	\begin{itemize}
		\item Email immernoch ein viel genutzes Kommunikationsmittel (Zahlen)
		\item Sehr viele Emails werden empfangen (Zahlen)
		\item Nutezende m�ssen das bearbeiten der Mails m�glichst effizient machen um andere Dinge zu schaffen
		\item Es gibt viele Email Clients, die versprechen die Email Kommunikation, schnellfer, effizienter und zukunftsorienteirter zu machen
		\item Emails Snoozen in vielen Clients ein Teil davon
		\item Obwohl viele die gro�en Anbieter nutzen (Zahlen), ist es gerade f�r diejenigen, die sich der Macht der gro�en entziehen wollen und z.b. auf Open Source wert legen notwendig offenen Stanndard f�r Email Management zu haben
		\item Den gibts auch, IMAP 
	\end{itemize}
	\item Gegenstad/Problemstelung/Ziele, Begr�ndung Wahl der Forschungsfrage indem Forschungsl�cke erl�utert wird \hl{Was ist der Gegenstand meiner Arbeit? Warum habe ich mich f�r diese Problemstelung entschieden und warum wichtig zu beantworten? Was ist das Ziel der Arbeit?}
	\begin{itemize}
		\item Ich bin selber fan von offenen Standards und der unabh�nigkeit von gro�en Unternehemen wie Google oder Microsoft
		\item deslab ist es wichtig zu schauen, inwieweit neue Features berits durch den existierenden ofenen standard umgsetzt werden k�nnen ud falls nicht, dass m�glichkeiten f�r erweitungen erarbeitet wrrden
		\item Das Problem, dass IMAP nur Grundfunktionalit�ten zur Verf�ugn stellt. Neue Features wie Snooze werden nicht unterst�tzt.
		\item Deshalb viele verschiedene Implementierungen dieses Features in den Cleints.
		\item untereinander nicht nutzbar (au�er leiches Unternehmen, sync �ber Account)
		\item vereinheitlichung w�re hier W�nschenswert
		\item In der Arbeit soll analysiert werden, wie die einzelnen Clients das Feature implementieren
		\item geschaut werden, inwieweit die Implementieurngen bereits interoperabel sind
		\item dann schauen, welche M�glichkeiten IMAP beitet um dieses Feature interoperabel umsetzen zu k�nnen
		\item Oder welche Erweiterungen geschaffen werden m�ssten
	\end{itemize}
	\item Gang der Untersuchung, Methodik skizzieren, Wie umfassend wurde die Literatur analysiert \hl{Wie l�sst sich die Methodik beschreiben, mit der ich meine zentrale Forschungsfrage analysiert bzw. beantowrtet habe?}
	\begin{itemize}
		\item Dazu wird zuerst analysiert, welche clients das feature wie implementiert haben
		\item dann geschaut welche m�gichekiten IMAP bietet das Feature zu implementieruen. bzw welche l�sungen missbraucht/adaptiert werden k�nnten
		\item was n�tig w�re als erweitung um das ganze zu standardtisieren
		\item das ganze basiert gr��tenteils auf aussagen der Client entwickler/Firmen oder testen und analysieren der clients
		\item zur standard analyse werden die rfcs zu IMAP benutzt, aber auch disussionen in (foren/Github)
		\item anhand der Nextcloud Mail App wird geschaut wie die analysierte interoperabilit�t implementiert werden k�nnte
	\end{itemize}
	\item Abgrenzung der Arbeit, Schwerpukt setzen, Bei bereits vorhandenen Arbeiten darauf hinweisen \hl{Muss ich meine Arbeit eingrenzen? Wie kann ich die Akzentuierung blausiebel begr�nden?}
	\begin{itemize}
		\item Schwerpunkt ist M�glichkeiten aufzeigen, wie das Snooze Feature �ber IMAP interoperabel implementiert werden k�nnte. Dabei sollen vorhandene Standards oder verbreitete Implementierungen vorzugsweise genutzt werden
		\item ...
	\end{itemize}
	\item Aufgabu der Arbeit, wie bin ich generell vorgegangen, um die Forscungsfrage zu beantworten, Damit ist im wesentlichen die Argumentationsline gemeint, die sich aus Ihrer Gliederung ergibt \hl{Wie habe ich meine ArbeitAufgebaut, wie werde ich vorgenen?}
	\begin{itemize}
		\item zuerst die grundlagen (Inbox Zero, IMAP, Email, interoperabel usw.) erkl�ren um das nachfolgende verstehen zu k�nnen
		\item dann erkl�ren, warum es einen interoperablen Ansatzt braucht und welche M�glichkeiten braucht man um Inbox Zero umzusetzen, wie sieht snoozen soll Snooze in Email clients aussehen?
		\item dann werden email cients analysiet und genasucht, wie die implementierung aussieht
		\item folgenden werden L�sungsan�tze beschrieben, die die snooze interoperabel m�glich machen k�nnten
		\item das ganze wird dann an der Nextcloud Mail app beschrieben
		\item der letzte schritt w�re die beschriebung einer Standardeerweiterung, was m�sste diese enthalten
	\end{itemize}
\end{itemize}



\chapter{AT: What is it, How it works}
\label{sec:HowItWorks}

\section{AT: Effizientes Mail Management}
\label{sec:InboxZero}

Viele stehen vor dem Problem, dass ihr E-Mail Posteingang \hl{(im folgenden INBOX genannt)} �berf�llt ist mit E-Mails, welche noch beantwortet werden m�ssen. Das E-Mail Postfach �ffnen f�hrt zu Stress, und es werden immer mehr E-Mails obwohl man sein bestes gibt diese zu bearbeiten. Es scheint so, als ob die F�higkeit mit der Masse an E-Mails klar zu kommen, entscheided ist �ber den Erfolg, den man vorallem im Job hat. Die Inbox wird als To-Do Liste, Kalender und Postfach in einem genutzt. \cite{googletechtalksInboxZero2007} Dieses Konzept f�hrt gezwungener Ma�en zu Problemen. \\
Es gibt Konzepte, mit denen dieses Problem angegangen werden kann. Die meisten basieren auf Getting Things Done, was eine Produktivit�tsstrategie ist, die das Ziel hat Aufgaben, Informationen, Probleme oder Projekte extern, also nicht im Kopf, zu verwalten. \cite{GettingThingsDone2023} David Allen hat diese Methode entwickelt und beschreibt in seinem Buch "Getting Things Done: The Art of Stress-Free Productivity", wie Sie funktioniert und anzuwenden ist. \hl{Buch ref?}\\
Auf der allgemeinen Getting Things Done Methode basieren andere Strategien, welche sich auf das E-Mail Management spezialisieren. Inbox Zero ist eine davon, welche von Merlin Mann erstmalig definiert wurde und nicht das Ziel hat keine E-Mails mehr in der Inbox zu haben, wie man es durch den Namen annehmen k�nnte, sondern eher darauf bezogen ist, kein Stress bei der Bearbeitung der Inbox zu haben. \cite{WhatInboxZero2019} Es geht nicht darum alle E-Mails in der Inbox zu beantworten sondern alle E-Mails zu bearbeiten. Das hei�t, dass man sich Aktionen definiert, welche man auf die zu bearbeitende E-Mail anwendet. Das kann das Beantworten der E-Mail sein aber genauso das L�schen oder Zur�ckstellen der E-Mail. Wichtig ist, dass man ohne viel �berlegen entscheidet, welche Aktion auszuf�hren ist und sich nach ausf�hren der n�chsten E-Mail widmet. \cite{googletechtalksInboxZero2007}
\hl{�hnlich ist die N days later Methode. (Nicht relevant?)}

\begin{itemize}
	\item Stress bei Inbox �ffnen, da viel zu viele Mails
	\item Verschiedene Methoden das zu �ndern, welche alle �hnlich sind und das gleiche Ziel verfolgen
	\begin{itemize}
		\item Getting Things Done
		\item Inbox Zero (productivity philosophy about dealing with a constant, life-long stream of email erstmalig von Merlin Mann erdacht) \cite{WhatInboxZero2019}
		\begin{itemize}
			\item Inbox ist nur ein �bergang, aber nicht das Ende bzw. der Endpunkt f�r eine E-Mail \cite{googletechtalksInboxZero2007}
		\end{itemize}
		\item To-Do Inbox
		\item (N-days later)
	\end{itemize}
\end{itemize}

\section{AT: E-Mail and IMAP}
\label{sec:EmailAndIMAP}

\begin{itemize}
	\item Wie wird Snooze definiert?
	\item Was ist IMAP und wie funktionierts
	\item Was genau hei�t interoperabel
	\item Was ist der Mailserver???
	\item INBOX/Posteingang Begriff erkl�ren Sysnonym
	\item Ordner/Mailbox Begriff erkl�ren Sysnonym
	\item Client erkl�ren?
	\item Tag/Flag/Label/Permaflag definieren und unterscheiden in Bezug auf IMAP oder unten bei thunderbird (zumindest permaflags)
	\item Auf Thread M�glichkeit eingehen, oder unten bei Mailspring
	\item RFCs hier erkl�ren? Was ist das und wie funktionierts
\end{itemize}


\chapter{AT: Problemstellung in Mail}
\label{sec:Problemstellung}

\begin{itemize}
	\item Wie kann kann Inbox Zero/GTD in Mail Clients aussehen, bzw. welche M�glichkeiten braucht man um diese Konzepte umsetzen zu k�nnen?
	\item Warum ist es n�tig/sinnvoll das �ber IMAP zu implementieren?
	\begin{itemize}
		\item Interoperabel
	\end{itemize}
\end{itemize}


\chapter{AT: Exsitierende Implementierungen}
\label{sec:ExistierendeImpl}

In diesem Kapitel werden verschiedenste Email Clients betrachtet und beschrieben ob und in welcher Form sie eine Snooze Funktionalit�t implementieren. Die Auswahl der betrachteten Clients basiert auf drei Schwerpunkten:\\
1. Popularit�t\\
Um den Status Quo analysieren zu k�nnen ist es wichtig verschiedene Clients zu betrachten. Allerdings gibt es nur eine Hand voll Email Clients, welche den Gro�teil des Marktes bei der Nutzung ausmachen. Auch wenn die Zahlen nicht ganz eindeutig sind, ist klar, dass die meisten Menschen die Clients von Gmail, Apple Mail und Outlook nutzen. Weitere Clients mit weniger Nutzenden, aber immernoch relevanten Zahlen sind Yahoo Mail, Web.de, GMX, Thunderbird und Samsung Mail. \cite{MostPopularEmail2022}  \cite{EmailClientMarketa} \cite{EMailNutzung2021Mailjet} \\
2. Open Source\\
Der zweite Schwerpunkt sind Open Source Clients. \\
3. Snooze Funktionalit�t\\
Der letzte Schwerpunkt sind Email Clients, welche die Snooze Funktionalit�t implementieren. Es gibt viele verschiedene Implementierungen f�r die Snooze Funktionalit�t. Um bewerten zu k�nnen, welche L�sungen gut und welche eher unbrauchbar sind, ist es gut einen Gesamt�berlick zu bekommen. Dabei k�nnen gerade Email Clients mit wenig Nutzenden eine sehr kluge Implementierung haben.

Die verschiedenen Clients werden in vier Gruppen unterteilt.\\
\textbf{Keine Implementierung:} Clients, welche keine Snooze Funktionalit�t haben.\\
\textbf{Lokale Implementierung:} Clients, welche eine Snooze Funktionalit�t haben, welche aber nicht f�r andere Clients benutzbar ist.\\
\textbf{\hl{Environment Implementierung}:} Client, welche die Snooze Funktionalit�t besitzen und diese Client�bergreifend zumindest in der eigenen Umgebung funktionieren.\\
\textbf{Interoperable Implementierungen:} Clients, welche die Snooze Funktionalit�t so implementiert haben, dass sie theoretisch von anderen Clients benutzt werden kann.


\begin{itemize}
	\item Kein IMAP
	\item Welche Annahmen wurden �ber propr�rit�re Software getroffen??
	\item wie wurde die snooze funktion getestet? Wurde traffic angeschaut? welcher mailserver (dovecot) wurde benutzt, welche extentions sind dort aktiviert? Wurde der metadaten ext gar nicht ausprobiert?!?
	\item Add Yahoo, Samsung und Windows Live Mail
\end{itemize}

\section{AT: Keine Implementierung}
\label{sec:NoImpl}

Bei der Analyse von bestehenden Implementierungen einer Snooze Funktionalit�t haben sich einige Email Clients herausgestellt, welche diese Funktionalit�t gar nicht implementiert haben. Darunter fallen der Android E-Mail Client K-9 \footnote{K-9: \url{https://k9mail.app/}}, OpenExchangeMail \footnote{Open-Xchange Email: \url{https://www.open-xchange.com/products/ox-cloud}}, Yahoo Mail \footnote{Yahoo Mail: \url{https://overview.mail.yahoo.com/}}, SnappyMail \footnote{SnappyMail: \url{https://snappymail.eu/}}, Rainloop \footnote{Rainloop: \url{https://www.rainloop.net/}}, Tutanota \footnote{Tutanota: \url{https://tutanota.com/}}, Posteo \footnote{Posteo: \url{https://posteo.de/de}}, ProtonMail \footnote{ProtonMail: \url{https://proton.me/mail}}, Cypht \footnote{Cypht: \url{https://cypht.org/}} oder Evolution \footnote{Evolution: \url{https://wiki.gnome.org/Apps/Evolution}}.


\section{AT: Lokale Implementierung}
\label{sec:LocalImpl}

Es gibt aber auch Email Clients, welche die Snooze Funktionalit�t implementiert haben aber diese nur in dem jeweiligen Client nutzbar ist. Die folgenden Clients fallen darunter:

\subsection{Mailbird}

Mailbird \footnote{Mailbird: \url{https://www.getmailbird.com/de/}} ist ein Desktop Email Client f�r Windows und Mac. Die Snooze Funktion ist nur in der Inbox verf�gbar und verschiebt die Email in in den Reiter \enquote{Zur�ckgestellt} bzw. \enquote{Snoozed}. �ber IMAP wird die Email in eine Mailbox \enquote{Snoozed} verschoben. 'Wenn die Zeit abl�uft, wird die E-Mail wieder ganz oben in Ihrem Posteingang angezeigt.' \cite{NtzenSnoozeFnktionIhrem2022} �ber IMAP ist dies nicht der Fall. Zus�tzlich wird die Email noch als ungelesen markiert.

\subsection{FairEmail} 
FairEmail\footnote{FairEmail: \url{https://email.faircode.eu/}} ist ein E-Mail Client f�r Android. Beim snoozen einer Mail wird diese f�r den Nutzenden ausgeblendet. F�r andere Email Clients ist dieses Ausblenden nicht sichtbar, da es lokal auf dem Ger�t passiert.\cite{FairEmailFAQMd} \\
Nachdem der definierte \hl{unsnooze} Zeitpunkt erreicht ist wird die Email wieder sichtbar und als ungelesen markiert. Au�erdem wird ein Uhr-Symbol angezeigt, welches darauf verweist, dass die Email \hl{gesnoozed} wurde.

\begin{itemize}
	\item 'snoozing is done locally on the device by hiding the message while it is snoozing. Unfortunately, it is not possible to hide messages on the email server too. Deleting messages from the server and restoring them later could result in losing messages when the device breaks.' \cite{FairEmailFAQMd}
\end{itemize}

\subsection{Thunderbird}
Thunderbird\footnote{Thunderbird: \url{https://www.thunderbird.net/de/}} hat keine Snooze Funktionalit�t. Allerdings gibt es ein Add-On mailmindr\footnote{mailmindr: \url{https://mailmindr.net/de/}}, welches lokales \hl{snoozen} in Thunderbird erm�glicht. Dabei werden die Emails in einen, vom Nutzenden definierten Ordner, verschoben. Beim \hl{unsnoozen} wird neben dem zur�ckschieben der Email noch eine Funktionalit�t ausgef�hrt, die vom Nutzenden definiert wird. Beispielsweise das Makieren als ungelesen der Email.\\
Es gibt in Thunderbird auch die M�glichkeit Emails zu taggen. Dazu gibt es f�nf vordefinierte Tags, wobei der f�nfte Tag \enquote{Later} hei�t. Intern wird der Tag als \enquote{\$label5} gef�hrt. Diese f�nf Label werden auch von anderen Email Clients wie z.B. der Nextcloud Mail App verwendet. In Thunderbird l�sst sich nach diesem Tag filtern, was auch eine M�glichkeit bietet Emails zu verstecken und zu einem anderen Zeitpunkt erneut anzuschauen.

\begin{itemize}
	\item Permaflags erkl�ren oder oben
\end{itemize}

\section{AT: Environment Implementierung}
\label{sec:EnvironmentImpl}

Clients, welche es auf verschiedenen Plattformen gibt, z.B. als App und Web Client haben meistens eine Syncronisation der Metadaten im Hintergrund �ber einen Nutzendenaccount. Dabei ist die Snooze Funktionalit�t in diesen Clienets meist interoperabel aber nicht �ber diese Umgebung hinaus. Die folgenden Clients bzw. Anbieter fallen darunter:

\subsection{GMail}
Gmail\footnote{Gmail: \url{https://www.google.com/intl/de/gmail/about/}} hat die Snooze-Funktionalit�t am 25.Juli 2018 eingef�hrt. \cite{NewGmailNow} Um zu verstehen, wie das \hl{snoozen} funktioniert, muss erst einmal erl�utert werden, wie Gmail generell Ordner/Mailboxen umsetzt. F�r die Organisation werden sogenannte Label angeboten, welche den Emails zugewiesen werden k�nnen. \cite{GmailMithilfeLabels} Diese Label sind dann vergleichbar mit Ordnern und k�nnen �ber das Men� einzeln gelistet werden. �ber IMAP werden die Labels als Ordner/Mailboxen Zur�ckgegeben. \cite{IMAPExtensionsGmail} Emails, welche beispielsweise das Label \enquote{Wichtig} zugeordnet bekommen, werden �ber IMAP als Email des Ordners/der Mailbox \enquote{Wichtig} gelisted.\\
Wenn eine Email in Gmail aus dem Posteingang \hl{gesnoozt} wird, wird das Label \enquote{Posteingang} entfernt und die Email wird in Gmail unter dem Reiter \enquote{Zur�ckgestellt} bzw. \enquote{Snoozed} gelisted. Falls die Email das Label \enquote{Posteingang} nicht mehr besitzt und ein anderes Label hat, wird dieses nicht entfernt. Durch das Entfernen des Labels \enquote{Posteingang} wird die Email �ber IMAP nicht mehr im Posteingang/Inbox angezeigt. Da Gmail den \hl{gesnoozten} Emails aber kein Label extra zuordnet, sind die Emails, welche aus dem Posteingang \hl{gesnoozt} werden, �ber IMAP nicht mehr auffindbar. Lediglich in einem Ordner, welcher alle auf dem Gmail Account existierende Email beinhaltet, werden die \hl{gesnoozten} Emails noch angezeigt. \\
Nachdem der definierte \hl{unsnooze} Zeitpunkt erreicht ist wird die Mail mit dem Label \enquote{Posteingang} versehen und wird damit auch f�r IMAP wieder sichtbar im Posteingang gelistet. Die Email wird nicht als ungelesen markiert aber in Gmail an den Anfang des Posteingangs sortiert.

Gmail gibt damit anderen Clients gar keine M�glichkeit �ber IMAP zu ermitteln, welche Emails gerade \hl{gesnoozt} sind.

\subsection{Exchange/Outlook}
Outlook\footnote{Outlook: \url{https://outlook.live.com/owa/}} erm�glicht es Emails aus der Inbox zu \hl{snoozen}. Die Emails werden dann in der Outlook App unter dem Reiter \enquote{Erneut erinnern} und in Outlook Webmail unter \enquote{Erneut erinnern aktiviert} gelistet. �ber IMAP werden die Emails in einen Ordner \enquote{Erneut erinnern aktiviert} verschoben, wobei sich der Name dieses Ordners je nach Spracheinstellungen �ndert. Wenn der \hl{unsnooze} Zeitpunkt erreicht ist, wird die Email zur�ck in die Inbox verschoben und als ungelesen markiert. In Outlook wird die Email dar�ber hinaus an den Anfang der Inbox sortiert.

\subsection{Apple Mail}
Apple Mail\footnote{Apple Mail: \url{https://support.apple.com/mail}} hat seit iOS 16 und macOS Ventura (13) eine \enquote{Remind me} Funktion.\cite{HowUseApple2022} In der Inbox befindliche Emails \cite{RemindMeOptions} lassen sich dabei mit einer Erinnerung versehen. Die Email wird dabei aber nur mit einem Uhr Symbol dargestellt und nach dem Zeitpunkt des \hl{unsnoozen} an den Anfang der Inbox sortiert. �ber IMAP ist die \enquote{Remind me} Funktion gar nicht zu erkennen.

\subsection{Mailspring}
Mailspring\footnote{Mailspring: \url{https://www.getmailspring.com/}} ist ein Desktop Email Client f�r Linux, Mac und Windows. Die Snooze Funktion ist sowohl in der Basic als auch in der Pro Version verf�gbar, wobei die Pro Version unbegrenztes \hl{snoozen} erlaubt.\cite{SnoozeEmailsHandle2021} \\
Emails, welche \hl{gesnoozt} werden, verschiebt Mailspring in einen Unterordner \enquote{Snoozed} des Ordners \enquote{Mailspring}. Beim \hl{unsnoozen} wird die Email zur�ck in die Inbox verschoben. Zus�tzlich wird eine neue Email generiert, welche �ber einen Thread der \hl{ungesnoozten} Email zugeordnet wird. Diese Email ist als ungelesen markiert und hat die aktuelle Zeit als Empfangszeit. Dadurch wird der Thread, und damit auch die \hl{ungesnoozte} Email, am Anfang der Inbox angezeigt, was sogar erm�glicht, dass auch �ber IMAP die Email wieder an den Anfang der Inbox verschoben wird.

\begin{itemize}
	\item Noch auf Thread Funktionalit�t eingehen oder oben erkl�ren
\end{itemize}

\subsection{Spark}
Spark\footnote{Spark: \url{https://sparkmailapp.com/de}} ist ein Email Client f�r Windows, Android, Mac und iOS, welcher auch \hl{snoozen} kann.\cite{EMailsFuerSpaeter} Die Snooze Funktionalit�t verschiebt die Email aus dem jeweiligen Ordner und zeigt sie unter dem Reiter \enquote{Snoozed} an. �ber IMAP wird die Email in einen Ordner \enquote{Later} verschoben. Beim \hl{unsnoozen} wird die Email zur�ck in die Inbox verschoben (egal ob sie aus der Inbox oder einem anderen Ordner heraus \hl{gesnoozed} wurde). Zus�tzlich wird die Email als ungelesen markiert und in den Spark Clients an den Anfang der Inbox sortiert. Syncronisiert werden die snooze Informationen �ber einen Spark Account, welcher automatisch beim einbinden des ersten Email Accounts angelegt wird. Wenn man bei einem anderen Spark Client den prim�ren Email Account hinzuf�gt, werden die snooze Informationen auch auf dem neuen Client �bernommen. \cite{SyncSparkMultipleb}

\subsection{Fastmail}
Fastmail\footnote{Fastmail: \url{https://www.fastmail.com/}} ist ein Email Client f�r Android, iOS und Web, welcher die Snooze Funktionalit�t als kostenpfichtiges Feature anbietet. Auch der IMAP Zugriff ist eine kostenpflichtige Funktionalit�t und l�sst sich im Gegensatz zu Snooze Funktionalit�t nicht mit einem kostenlosen Account testen. Laut Fastmail ist die Snooze Funktionalit�t auf Fastmail beschr�nkt und andere Clients, welche z.B. �ber IMAP auf den Account zugreifen, haben keinen Zugriff auf diese Funktionalit�t.\cite{SnoozingMail}
Wenn eine Email in Fastmail \hl{gesnoozt} wird, dann wird sie aus dem Ordner in dem sie sich befindet ausgeblendet und unter dem Reiter \enquote{Snoozed} gelistet. Beim \hl{unsnoozen} wird die Email wieder in dem Ursprungsordner angezeigt und an den Anfang des Ordners sortiert. Der gelesen Status wird nicht ge�ndert.

\begin{itemize}
	\item (Behauptung Fastmail benutzt IMAP nutzt draft-ietf-sieve-imap-sieve-09 kombiniert mit draft-ietf-extra-sieve-snooze)
	\item \hl{Hier evt auf JMAP eingehen?}
\end{itemize}

\subsection{Zoho Mail}
Zoho Mail\footnote{Zoho Mail: \url{https://www.zoho.com/de/mail/}} bietet die Snooze Funktionalit�t. \cite{SnoozeEmailsLater} Aber auch hier ist der IMAP Zugriff ein kostenpflichtiges Feature. \cite{ZohoMailFlexibles} In Zoho Mail wird beim \hl{snoozen} die Email unter dem Reiter \enquote{Snoozed} gelistet und aus dem ehemaligen Ordner entfernt. Beim \hl{unsnoozen} wird die Email in die Inbox verschoben (egal, aus welchem Ordner heraus sie \hl{gesnoozt} wurde) und der ungelesen Status wird nicht ge�ndert. Auch wird die Email nicht an den Anfang der Inbox sortiert.


\section{AT: Interoperable Implementierungen}
\label{sec:InteropImpl}

Im folgenden werden Implementierungen in Email Clients betrachtet, welche einen interoperablen Ansatz verfolgen und somit theoretisch von allen Email Clients benutzt werden k�nnen.

\subsection{Roundcube}
Roundcube\footnote{Roundcube: \url{https://roundcube.net/}} ist ein selbstgehosteter Webmailer und kann von sich aus nicht snoozen. Es gibt allerdings ein Plugin roundcube-snooze\footnote{roundcube-snooze: \url{https://github.com/toptensoftware/roundcube-snooze}}, mit dem es m�glich ist Emails zu \hl{snoozen}. Beim \hl{snoozen} wird dabei der Header der Email bearbeitet. Da Emails aber nicht bearbeitbar sind wird die Email dubliziert und der Header wird bearbeitet, bevor die neue Email gespeichert wird. Die alte Email wird gel�scht. Beim bearbeiten wird ein \hl{Header-Tag} \textit{X-Snoozed} hinzugef�gt, welcher den Zeitpunkt des Snooze als auch den zuk�nftigen Zeitpunkt des \hl{Unsnooze} enth�lt. Au�erdem wird der Name der Mailox, aus dem die Email \hl{genoozt} wurde und der aktuelle Status gespeichert.\cite{robinsonSnoozeRoundcube2023} Im folgenden ist der Header-Tag einer Beispiel-Email zu sehen, welche bereits \hl{unsnoozt} ist.

\begin{verbatim}
  X-Snoozed: at Fri, 13 May 2022 10:16:56 +1000;
    until Fri, 13 May 2022 10:17:00 +1000;
    from INBOX;
    woken
\end{verbatim}

Zus�tzlich wird die Email in eine Snoozed Mailbox verschoben.\\
Ein Aufwachskript, welches auf dem Webserver des Webmailers installiert wird, schiebt die Email nachdem der \hl{unsnooze} Zeitpunkt erreicht ist zur�ckn und das Keyword \textit{woken} wird im Header-Tag \textit{X-Snoozed} gesetzt.\cite{robinsonSnoozeRoundcube2023}\\
Dadurch, dass der Header der Email bearbeitet wird, k�nnen alle anderen Email Clients diese Information sehen und dementsprechend die Informationen theoretisch darstellen.\\
Die entwickelnde Person macht aber klar, dass es sich um eine sehr experimentelle Implementierung handelt, welche nicht ausgibig getestet wurde und nur auf DoveCot\footnote{DoveCot: \url{https://www.dovecot.org/}} ausgerichtet ist.\cite{robinsonSnoozeRoundcube2023}


\subsection{eM Client}
Der eM Client\footnote{eM Client: \url{https://de.emclient.com/}} ist ein Email Client f�r Windows und Mac. Mobile Clients f�r Android und iOS sind aber bereits in der Beta verf�gbar. \cite{EMClientBeta} \\
Die Snooze Funktionalit�t ist ein kostenpflichtiges Feature in der Pro Version. \cite{Snooze} Allerdings l�sst sich in der Android App auch ohne Pro \hl{snoozen}. Beim \hl{snoozen} wird der Email ein \hl{Label} zugeweisen,\hl{(Permaflag, ref auf oben)} welches mit \enquote{\$SnoozedUntil\_} beginnt und von dem \hl{Unsnooze} Zeitpunkt als Unix Timestamp gefolgt wird. Daraus entsteht beispielsweise das folgende Label: \enquote{\$SnoozedUntil\_162597600}\\
In den eM Clients werden die Emails mit solch einem Label, bei dem der Timestamp in der Zukunft liegt, in der Inbox ausgeblendet. Die Email kann dann unter dem Reiter \enquote{Snoozed} gefunden werden. Nach dem \hl{unsnoozen} wird die Email wieder sichtbar.\\
Da das Label durch andere Email Clients abgerufen werden kann, ist es theoretisch m�glich, dass auch andere Clients die Email \hl{gesnoozt} darstellen.


\section{AT: Fazit}
\label{sec:FazitExistingSolution}

Wie man sehen kann, gibt es viele verschiedene Implementierungen einer Snooze Funktionalit�t. Teilweise sehr unterschiedlich aber oft auch mit gleichen Ans�tzen. Generell f�llt auf, dass oft versucht wird auf lokale Implementierungen zu setzen. Viele Clients legen eine extra Mailbox an um die gesnootzen Emails dorthin zu verschieben. Das ist in gewisser Weise eine M�glichkeit f�r andere Clients zumindest zu erkennen, welche Emails gesnoozt sind und welche nicht. Allerdings sind die Namen dieser Mailboxen nicht einheitlich und unterscheiden sich sogar teilweise in der eigenen Umgebung, wie es bei Outlook der Fall ist. Auch das Verhalten der Clients nach dem unsnoozen unterscheidet sich stark. Vom reinen wieder einblenden bis als ungelesen an den Anfang der Mailbox sortieren und mit einem besonderen Symol makieren ist alles dabei.
Nichts destotrotz gibt es auch L�sungsans�tze, welche versuchen eine interoperable Implementierung zu nutzen. Auf diese und andere interoperable Ans�tze wird im n�chsten Kapitel n�her eingegangen.

\chapter{AT: Diskussion L�sungsans�tze}
\label{sec:DisskussionLoesung}

Das folgende Kapitel widmet sich, nach der Analyse existierender L�sungen im vorherigen Kapitel, dem schaffen/beleuchten/analysieren interoperablen L�sungen f�r eine Snooze Funktion. Dabei werden Ans�tze von den bereits existierenden Implementierugen noch einmal aufgegriffen und neue/andere Ans�tze beleuchtet. Das Ziel ist, dass die Ans�tze/L�sungen nicht dazu f�hren, dass andere Clients, welche diesen Ansatz/L�sung nicht implemntiert haben verwirrt werden \hl{FORMULIERUNG} und wie bei Gmail z.B. gesnoozte emails nicht mehr aufzufinden sind. Im besten Fall ist auch f�r Clients ohne die Implementierung eines Ansatztes erkennbar, dass eine Email gesnoozt ist. \\
Der Fokus liegt auf der Verwendung bereits existierender Funktionalit�ten im IMAP Protokoll oder Erweiterungen dieser durch sogenannte Extentions. Aber auch das Ausnutzen der Strucktur einer Email ist ein bevorzugter Ansatz. Grund daf�r ist, dass je n�her man am Standard bleibt oder auf bereits etablierte Erweiterungen setzt, umso wahrscheinlicher ist am Ende auch, dass die eigene Implementierung eines Ansatzen Verbreitung findet.

\section{AT: Was muss m�glich sein}
\label{sec:WasMussMoeglich}

Um die L�sungsans�tze bewerten zu k�nnen ist es n�tig nochmal klar zu formulieren, welche Aufgaben eine Snooze Funktion erf�llen muss und welche Funktionen dar�ber hinaus w�nschenswert w�ren, damit man von einem interoperablen Ansatz eines Snooze Features sprechen kann.\\
Der erste erforderliche Aufgabe ist, dass erkennbar sein muss, welche Email gesnoozt ist und welche nicht. Auch muss ermittelbar sein, bis wann die Email gesnoozt ist, also wann sie wieder unsnoozt werden soll. Aus welcher Mailbox eine Email gesnoozt wurde, ist wichtig, wenn die Email beim Snoozen in eine andere Mailbox verschoben wird. Wie in der Analyse der existierenden Implementierungen beschrieben gibt es f�r dieses Verhalten keine Einheitlichkeit. Manche Clients verschieben in die Mailbox zur�ck, aus der gesnoozt wurde, andere konsequent in die INBOX. F�r den Fall, dass diese Information nicht ermittelbar ist kann die Email aber immer in die Inbox verschoben werden, weshalb diese Information nicht unbedingt vorhanden sein muss. Auch die Information, seit wann eine Email gesnoozt ist, ist f�r die Funktionalit�t der Snooze Funktion nicht entscheident, k�nnen durch den Client bei Verf�gbarkeit aber dargestellt werden.\\
Zum Schluss muss sichergestellt werden, dass beim unsnooze Prozess keine Inkonsitenten Daten enstehen, weil z.B. mehrere Clients gleichzeitig eine Email unsnoozen.

\section{AT: Welche M�glichkeiten gibts}
\label{sec:WelcheMoegGibts}

\subsection{AT: Keywords}

Der erste L�sungsansatz, der beleutet wird ist die Verwendung von \hl{Keywords}. Dabei werden die durch IMAP definierten \hl{Permaflags} genutzt. Wie bereits erl�utert k�nnen diese von den Clients erstellt und den einzelenen Email zugeordnet werden. Es ist zwar m�glich, dass der Mailserver das erstellen eigener \hl{Permaflags} nicht gestattet, was aber nur sehr selten der Fall ist.

Angeleht an die Implementierung in Thunderbird, wo es ein \hl{Label} \enquote{Later} gibt, k�nnte hier ein \hl{Label} \enquote{Snooozed} eingef�hrt werden. Dieses \hl{Label} w�rde dann die Aufgabe erf�llen Emails zu identifizieren, welche \hl{gesnozt} sind. Es ist auch m�glich direkt das \hl{Label} \enquote{Later} von Thunderbird zu adaptieren f�r diesen Zweck. Da Later im Fall von Thunderbird eine sehr �hnliche Bedeutung hat wie ein Snooze hat und intern als \enquote{\$label5} implementiert ist, k�nnten dieses \hl{Label} auch \enquote{Snoozed} hei�en und somit direkt in mehreren Clients funktionieren. Zu beachten ist hier aber, dass dieses \enquote{Label} nur die Information repr�sentiert, welche Mail \hl{gesnoozt} ist, nicht aber f�r wie lange. Deshlab kann dieser Ansatz h�chstens ein Teil einer interoperablen L�sung sein.

Ein anderer Ansatz basierend auf der Verwendung von \hl{Keywords}, ist die Implementierung des eM Clients. Das \hl{Label}, welches dort beim \hl{snoozen} gesetzt wird besteht aus der ersten H�lte \enquote{\$SnoozedUntil\_} und der zweiten H�lte, welche ein UNIX Timestamp ist, z.B. \enquote{162597600}. Dabei w�rde die erste H�lte die Aufgabe erf�llen, zu identifizeiren, welche Email \hl{gesnoozt} ist und die zweite H�lfte die Zeit definieren, wann die Email \hl{unsnoozt} werden soll. Das Problem dieses Ansatzen liegt im Design des \hl{Labels}. Da der UNIX Timestamp Teil des \hl{Labels} ist, ist jedes \hl{Label} unterschiedlich. Ausgenommen der Fall, dass zwei oder mehr Emails gleichzeitig \hl{gesnoozt} gesnoozt werden. Dies f�rt dazu, dass bei jedem \hl{snoozen} ein neues \hl{Label} erstellt wird, welches dann durch den Mailserver gespeichert und bei jeder \enquote{SELECT} Anfrage zur�ckgegeben wird. Zwar definiert der IMAP Standard keine maximale Anzahl an \hl{Labeln}, welche durch die Clients erstellt werden d�rfen, allerdings kann der Mailserver das Erstellen neuer \hl{Label} unterbinden, falls die Anzahl eine interne grenze �berschreitet, indem die \hl{PERMAFLAG} Antwort kein \textbackslash* enth�lt. \cite{melnikovInternetMessageAccess2021} Au�erdem nutzen manche Clients wie z.B. FairEmail die \hl{PERMAFLAG} Antwort um eine Liste an m�glichen \hl{Labels} zu generieren, die einer Email zugewiesen werden k�nnen. Ein Heruasfiltern dieser \hl{Label} ist in diesem Fall wahrscheinlich ein annehmbarer Aufwand, den die Clients betreiben m�ssten. Aber auch generell entsteht durch diesen Ansatz eine sich schnell aufsummierende Masse an Daten, sowohl auf dem Mailserver als auch im Datenverkehr zwischen Client und Server, welche nur bis zu einem bestimmten Zeitpunkt gebraucht werden. Eine M�glichkeit \hl{Labels} zu l�schen oder ein mechanismus, der daf�r sorgt, dass ncht mehr verwendete \hl{Labels} Serverseitig gel�scht werden, w�re f�r diesen Ansatz w�nschenswert und notwendig.

\begin{itemize}
	\item Was sind die IMAP Keywords??!? Wie unterscheiden sich Flags/Permaflags und der die anderen Begriffe (eher oben erkl�ren, da bereits verwendet)

\end{itemize}

\subsection{AT: Header bearbeiten}

\begin{itemize}
	\item evt header n�her erkl�ren (oder oben)
\end{itemize}

Der n�chste L�sungsansatz, ist das Bearbeiten des Email Headers wie es das roundcube plugin roundcube-snooze umsetzt. Dabei wird der Email, welche \hl{gesnoozt} wird ein neues \hl{Header Feld} \textit{X-Snoozed} hinzugef�gt. Im folgenden ist noch einmal das bereits verwendete Beipiel aufgef�hrt: 

\begin{verbatim}
  X-Snoozed: at Fri, 13 May 2022 10:16:56 +1000;
    until Fri, 13 May 2022 10:17:00 +1000;
    from INBOX;
    woken
\end{verbatim}

Das neue Header Feld enth�t sowohl den genau Zeitpunt des \hl{snooze} (at) als auch des \hl{unsnooze} (until) in Form von einem in RFC2822\footnote{RFC2822 - Internet Message Format: \url{https://datatracker.ietf.org/doc/html/rfc2822}} definierten Timestamps. Zus�tzlich wird der Name der Mailbox gespeichert, aus dem die Email gesnoozt wurde (from). Das Schl�sselwort \enquote{woken} wird dem Header Feld angeh�ngt, wenn die Email wieder \hl{unsnoozt} wurde. Mit diesem Ansatz und den gespeicherten Informationen lassen sich alle definierten Anforderungen an ein Snooze Feature realisieren.

Da dieser L�sungsansatz die Header einer Email bearbeitet, sind auch noch weitere Einfl�sse zu beachten. Es k�nnte die Frage aufkommen, in wie weit durch das Bearbeiten der Header Verschl�sselung und Signierung durch S/MIME beieinflusst werden k�nnen. Diesbez�glich ist bei richitger Implementierung allerdings keine Gefahr zu bef�hrchten. S/MIME kann durch das einschlie�en bestimmter Header in einen sogenannten \enquote{message/rfc822 wrapper}, auch nicht Inhaltsbezogene Header wie z.B. den \enquote{To} Header sch�tzen. \cite{schaadSecureMultipurposeInternet2019} \cite{resnickInternetMessageFormat2001} Allerdings sind diese explizit auszuw�hlen. Es gibt �u�ere Header, welche f�r S/MIME nicht relevant sind. Wenn der \textit{X-Snoozed} Header als einer dieser �u�eren hinzugef�gt wird, ergibt sich somit kein Einfluss auf S/MIME.\\
Auch Erweiterungen, wie RFC7508\footnote{RFC7508 - Securing Header Fields with S/MIME \url{https://datatracker.ietf.org/doc/rfc7508/}} basieren auf der Auswahl bestimmter Header, welche gesch�tzt werden und es erlauben, Header hinzuzuf�gen ohne die Sicherheitsmechanismen \hl{kaputt} zu machen. \cite{cailleuxSecuringHeaderFields2015}

Ein Problem dieses Ansatztes ist, dass die Header einer Email eig. nicht durch den empfangenden Client bearbeitet werden. \hl{Dazu gibts ja IMAP und Erweiterungen, um die Nachricht an sich unber�hrt zu lassen... Hier irgendwie noch begr�nden}

\subsection{AT: Metadaten Speichern}

Ein weiterer L�sungsansatz ist die Verwendung der RFCs 5257\footnote{RFC5257 - Internet Message Access Protocol - ANNOTATE Extension \url{https://datatracker.ietf.org/doc/html/rfc5257}} (Internet Message Access Protocol - ANNOTATE Extension) und 5464\footnote{RFC5464 - The IMAP METADATA Extension \url{https://datatracker.ietf.org/doc/html/rfc5464}} (The IMAP METADATA Extension). Beide Erweiterungen funktionieren sehr �hnlich. Die Annotate Extention erm�glicht zu jeder Email seperat Informationen abzuspeichern, die Metadata Extention dagegen zu Mailboxen und auf Server Ebene. Die Annotate Extention ist eine Experimentelle Erweiterung von 2008.\cite{dabooInternetMessageAccess2008} Aufgrund dessen f�llt die Implementierung in Mailservern eher gering aus.\cite{boschSpecsUnofficialIMAP} Gerade DoveCot als verbreitetster IMAP Mailserver \hl{BELEG} unterst�tzt diese Erweiterung nicht. Die Metadata Erweiterung dagegen ist Teil des Standards und wird von DoveCot implementiert.\cite{boschSpecsUnofficialIMAP} Es wird im folgenden der L�sungsansatz anhand der Metadata Erweiterung erkl�rt, obwohl die Annotate Erweiterung eine bessere Grundlage f�r das Snooze Feature bietet, da Informationen direkt auf Email Ebene abgespeichert werden k�nnen. F�r die Metadaten Erweiterung spricht in diesem Fall aber die bereits existierende Standardisierung und somit vorhandene Implementierung in vielen IMAP Mailservern.

Die Metadata Erweiterung f�hrt neue IMAP Befehle ein f�r das setzen und abrufen der Metadaten. Um auf Informationen auf Mailbox Ebene zuzugreifen ist der Mailboxname als erstes Argument mitzugeben. Um auf die Informationen auf Serverebene zuzugreifen wird ein leerer String als Argument mitgegeben. Der Speicher pro Mailbox bzw. f�r den Server ist dann Hirarchisch aufgebaut und besteht aus Eintragspaaren bestehend aus Name und Wert. Der Name des Eintrages beginnt mit einem \enquote{/} und wird von Worten gefolgt, wie wiederum mit \enquote{/} getrennt sind. Dieser Aufgabau bildet dann die Hirarchie der Eintr�ge.\\
Die oberste Hirachrchie-Ebene besteht aus zwei M�glichkeiten: \enquote{/private} und \enquote{/shared}. Unter private sind Eintr�ge zu findne, welche nur durch den eigeloggten Nutzenden abgerufen werden k�nnen. Unter shared befinden sich Eintr�ge, welche f�r alle Nutzenden, die Zugriff auf diese Mailbox oder den Server haben abrufbar sind. Die n�chste Hirarchiebenede besteht aus vordefinierten M�glichkeiten wie z.B. \enquote{/private/comment} oder \enquote{/shared/comment} f�r eine Notiz zur Mailbox. An dieser Stelle k�nnte man ansetzten und eine Eintrag \enquote{/private/snoozed} f�r eine spezielle Snoozed Mailbox erzeugen. In diese Mailbox werden die \hl{gesnoozten} Emails verschoben und als Wert des Eintrages wird ein JSON Objekt gespeichert, welches zu allen in der Mailbox befindlichen Emails die notwendigen Snooze Informationen abspeichert.

\begin{itemize}
	\item Evt noch drauf eingehen, dass Wert von Eintrag mehrzeitlig sein kann und string NIL oder binary data sein kann
	\item Evt auf min von 10 Eintr�ge pro Mailbox m�ssen erlaubt sein, dar�er hinaus vom server einschrenkbar, eingehen
\end{itemize}


\subsection{AT: Referenze Email}

Ein weiterer Ansatz ist das Speichern der Snooze Informationen in einer extra Email �hnlich zur Implementierung, die Mailspring verwendet um Emails an den Anfang der Inbox zu sortieren. Beim snoozen w�rde dabei eine neue Email erzeugt werden, welche \hl{per Referenz} der eigentlichen Email angeh�ngt wird. �ber ein Attribut dieser neuen Email, welches defineirt ist kann dann die Email im Client ausgeblendet, sodass nur die eigentliche \hl{gesnoozte} Email sichtbar ist. F�r dieses Attribut gibt es mehrere M�glichkeiten. Beispielsweise kann eine \hl{Keyword} \enquote{Hidden} eingef�hrt werden, welches den Clients mitteilt, dass die Email mit diesem \hl{Keyword} unsichtbar sein soll. Es kann auch ein Header in der Email einen bestimmten Wert haben. Beispielsweise k�nnte man den Subject Header auf \enquote{SNOOZEDMETADATA} setzen. Wichtig ist, dass dieses Attribut einheitlich und eindeutig definiert ist. Im Body der Email k�nnten dann die Informationen zum Snooze der eigentlichen Email in Form z.B. eines JSNON Objekts stehen.

F�r ander Clients, die diesen Ansatz nicht implementieren, ensteht wie auch bei Mailspring eine extra Email, die eigentlich kein relevanten f�r den Nutzenden hat. Viele Email Clients bieten implemntieren zwar die Thread Funktionalit�t und w�rden die extra Emial nur bei Auswahl der eigentlichen Email anzeigen, trotzdem kann dies zu Un�bersichtlichkeit und Verwirrung f�hren.

\subsection{AT: Dedicated Ordner}

Alle bis jetzt beschriebenen L�sungsans�tze, mit ausnahme der Benutzung der Metadata Erweiterung, erf�llen die Aufgabe, erkennbar zu machen, welche Email gesnoozt ist, ohne eine extra Mailbox anzulegen. Das hei�t aber auch, dass Clients, welche keine Implementierung einer dieser Ans�tze verfolgen, zu Verwirrung und Un�bersichtlichkeit f�hren k�nnen, wenn mehrere Clients benutzt werden. \cite{SnoozedEmailsStay2022} Es ist deshalb sinnvoll eine extra Mailbox anzulegen, welche \hl{gesnoozte} Emails aufbewahrt, bis sie wieder \hl{unsnoozed} werden. Dieses Verhalten gibt es bereits in vielen Clients, welche ein Snooze Feature implementiert haben. Allerdings gibt es in der Namensgebung dieser Mailbox, wie bereits erw�hnt, viele Unterschiede. Ein einheitlicher Name wie z.B. \enquote{Snoozed} w�re hier sinnvoll.

F�r eine interoperable L�sung kann sich an der Idee der \enquote{Special-Use Mailbox} orientiert werden, welche im  RFC6154\footnote{RFC6154 - IMAP LIST Extension for Special-Use Mailboxes: \url{https://datatracker.ietf.org/doc/rfc6154/}} definert wird und bereits Teil des offiziellen IMAP Standards in der neusten Version (4rev2\footnote{RFC9051 - Internet Message Access Protocol (IMAP) - Version 4rev2: \url{https://datatracker.ietf.org/doc/html/rfc9051}}) ist. Dabei gibt der Server beim Listen der Mailboxen bestimmte Attribute zur�ck, welche n�her beschreiben, was in der Mailbox zu finden ist oder was dort hinein geh�rt. \cite{melnikovInternetMessageAccess2021} Es gibt definierte fest definierte Attribute, wie z.B. \enquote{\textbackslash{Archive}} oder \enquote{\textbackslash{Drafts}}, welche eine Mailbox f�r archivierte Nachrichten und eine Mailbox f�r Entw�rfe spezifiziert. An dieser Stelle k�nnte eine Erweiterung des Standards um ein Attribut \enquote{\textbackslash{Snoozed}} angestrebt werden, welches dann f�r alle Clients klar definiert, welche Mailbox die Aufgabe hat \hl{gesnoozte} Emails aufzubewahren.

\subsection{AT: Savedate}

Eine Aufgabe, die f�r eine Snooze Funktionali�t nicht entschiedend ist, aber zu einem besseren Nutzendenerlebnis breitr�gt ist das sortieren an den Anfang nach dem \hl{unsnoozen}. Daf�r beitet sich die IMAP Erweiterung RFC8514\footnote{RFC8514 - Internet Message Access Protocol (IMAP) - SAVEDATE Extension: \url{https://datatracker.ietf.org/doc/html/rfc8514/}} (Internet Message Access Protocol (IMAP) - SAVEDATE Extension) an. Diese Erweiterung definiert ein neues Email Attribut \enquote{save date}. Dieser Zeitpunkt ist ist der Zeitpunkt, an dem die Email in die Mailbox gekommen ist, in der sie sich gerade befindet. Sobald eine Email also von einer Inbox in eine andere verschoben wird, �ndert sich dieses Attribut auf die Aktuelle Zeit. Abgerufen werden kann dieses Attribut �ber den \textit{FETCH} Befehl.\\
Wenn man davon ausgeht, dass eine Email beim \hl{snoozen} in eine andere Mailbox verschoben wird und beim \hl{unsnoozen} wieder zur�ckgeschoben wird, kann man die Emails der Mailbox nach diesem Attribut sortieren und w�rde Emails, welche gerade \hl{gewaked} sind am Anfang der Mailbox angezeigt bekommen.\\
Leider ist diese Erweiterung im Moment nur in DoveCot ab der Version 2.3.11 genutzt werden. \cite{boschSpecsUnofficialIMAP}

\subsection{AT: Wer �bernimmt das verschieben}

Alle beschriebenen L�sungsans�tze bieten teilweise oder vollumf�nglich die Grundlage f�r eine Snooze Funktionalit�t. Allerdings ist der eigentliche Prozess des \hl{snoozen} und \hl{unsnoozen} noch nicht geregelt. Dies kann entweder durch den Server oder durch den Client passieren. F�r den Fall, dass eine Clientseitige Implementierung angestrebt wird, ist es notwendig, dass alle Clients schnelle Updates �ber eine Ver�nderung auf dem Server bekommen, da das \hl{unsnoozen} theoretisch von vielen Clients gleichzeitig ausgef�hrt werden kann.\\
Hierf�hr bietet sich der RFC7162\footnote{RFC7162 -IMAP Extensions: Quick Flag Changes Resynchronization (CONDSTORE) and Quick Mailbox Resynchronization (QRESYNC): \url{hhttps://datatracker.ietf.org/doc/html/rfc7162}} an. Diese Erweiterung ist Teil des Standards und weit verbreitet implementiert.\cite{boschSpecsUnofficialIMAP} Die \enquote{Conditional Store} Funktionalit�t erm�glicht es mehrere �nderungen innerhalb einer Mailbox zu koordinieren und sicherzustellen, dass immer nur ein Client schreibzugriff hat. Zus�tzlich wird ein Mechanismus definiert, welcher es erm�glicht Mailbox �nderungen in Bezug auf Flag und Metadaten schnell zu synconisieren.\cite{melnikovIMAPExtensionsQuick2014}

\section{Fazit}
\label{sec:LoesungFazit}

\begin{itemize}
	\item 
\end{itemize}


\begin{itemize}
	\item (JMAP (nicht IMAP aber, sollte erw�hnt werden), vlt aber eher bei sparkmail.)
\end{itemize}

\section{AT: Nextcloud Mail App}
\label{sec:LoesungNextcloud}

\begin{itemize}
	\item Welcher L�sungsansatz passt am besten in die App
	\item horde (was kann, was nicht)
	\item 
\end{itemize}


\chapter{AT: Vorschlag Standarderweiterung}
\label{sec:VorschlagStandard}

\begin{itemize}
	\item Was ist RFC?
	\item Wie k�nnte ein RFC f�r eine Snooze aussehen
	\begin{itemize}
		\item Flags erweitern um Snoozed
		\item Special use mailboxes erweitern um Snoozed oder zumindest einheitlicher Name
	\end{itemize}
\end{itemize}
					% hier steht der eigentliche Text der Arbeit

\pagenumbering{Roman}
\appendix	
												
\bibliography{Literatur/Literatur} 	% Literaturverzeichnis

% Empfehlung: Verzeichnisse an das Ende der Arbeit (PKW)
% Ein Glossar ist eher untypisch fuer Graduierungsarbeiten
% Abkuerzungsverzeichnis an das Ende der Arbeit; Abkuerzungen und Begriffe muessen in jedem
% Fall im Haupttext der Arbeit beim ersten Vorkommen definiert werden

%% \printglossary[title=Glossar] 			% Glossar Einträge in Header/Glossar.tex vornehmen
%%\printglossary[type=\acronymtype,title=Abk\"urzungsverzeichnis]	% Abkürzungsverzeichnis Einträge in Header/Abkuerzungen vornehmen

\listoffigures											% Abbildungsverzeichnis
\listoftables												% Tabellenverzeichnis


\pagestyle{empty}

\chapter*{Selbst�ndigkeitserkl�rung}
\label{sec:Selbst�ndigkeitserkl�rung}

Ich erkl�re hiermit, dass ich die vorliegende Graduierungsarbeit ohne Hilfe Dritter und ohne Benutzung anderer als der angegebenen Quellen und Hilfsmittel verfasst habe. Alle den benutzten Quellen w�rtlich oder sinngem�� entnommene Stellen sind als solche einzeln kenntlich gemacht.

Diese Arbeit ist bislang keiner anderen Pr�fungsbeh�rde vorgelegt und auch nicht ver�ffentlicht worden.

Ich bin mir bewusst, dass eine falsche Erkl�rung rechtliche Folgen haben wird.\\[1.5cm]


\begin{minipage}{0.7\textwidth}
\ort, \today \hfill  Unterschrift
\end{minipage}





\end{document}